\section{Fitted Q-Iteration y redes Q profundas}

\subsection{De las tablas a la aproximación de funciones}

Cuando el espacio de estados es continuo o muy grande, no es posible mantener una tabla de valores Q. La solución consiste en parametrizar la función Q con un aproximador de funciones, por ejemplo, una red neuronal.

\begin{importantbox}{Fitted Q-Iteration (FQI)}
FQI convierte Q-Learning en un problema de aprendizaje supervisado:
\begin{enumerate}
    \item Con la red Q actual se calcula el TD target: $y_i = r_i + \gamma \max_{a'} Q_{\theta^-}(s_i', a')$
    \item Se toma $(s_i, a_i)$ como entrada y $y_i$ como etiqueta, y se entrena la red Q:
    $$\min_\theta \sum_i \left(Q_\theta(s_i, a_i) - y_i\right)^2$$
    \item Se actualiza la red objetivo $\theta^- \leftarrow \theta$
    \item Se repite
\end{enumerate}
\end{importantbox}

\subsection{Técnicas fundamentales de DQN}

\begin{importantbox}{Innovaciones de DQN (Deep Q-Network)}
DQN introdujo dos técnicas fundamentales para que el Q-Learning profundo funcione de manera estable:
\begin{enumerate}
    \item \textbf{Experience Replay}: la experiencia histórica se almacena en un replay buffer y, durante el entrenamiento, se muestrean mini-batches al azar. Esto rompe la correlación temporal de los datos y mejora el aprovechamiento de las muestras.
    \item \textbf{Target Network}: se usa una red objetivo $Q_{\theta^-}$, actualizada de forma periódica, para calcular el TD target. Esto reduce la inestabilidad que provoca el problema del «objetivo móvil».
\end{enumerate}
\end{importantbox}

\begin{warningbox}{Inestabilidad de Q-Learning con aproximación de funciones}
A diferencia del caso tabular, Q-Learning con aproximación de funciones \textbf{no garantiza la convergencia}. La propiedad de contracción de la actualización de Bellman puede dejar de cumplirse con aproximación de funciones, lo que provoca:
\begin{itemize}
    \item Divergencia de los valores Q (acumulación del overestimation bias)
    \item Entrenamiento inestable (oscillation o divergence)
    \item Sensibilidad a los hiperparámetros
\end{itemize}
Por eso se necesitan técnicas de estabilización como la target network y la experience replay.
\end{warningbox}

\subsection{Mejoras de DQN}

\begin{itemize}
    \item \textbf{Double DQN}: la red actual selecciona la acción y la red objetivo evalúa su valor, lo que reduce la overestimation
    \item \textbf{Dueling DQN}: descompone la función Q en $V(s) + A(s, a)$ y estima por separado el valor del estado y la ventaja
    \item \textbf{Prioritized Experience Replay}: muestrea con prioridad según la magnitud del TD error, para concentrarse en la experiencia de la que más se puede aprender
\end{itemize}

\subsection{Resumen de la sección}

El paso central del Q-Learning tabular a DQN es la aproximación de funciones. La experience replay y la target network son las técnicas fundamentales para que el entrenamiento sea estable.

